\addtocontents{toc}{\protect\Needspace{3\baselineskip}}
\originalchapter{第一次测验复习：阅读、相对论与物理图像补题}
本章补入 8.286 Fall 2013 第一次测验复习卷中此前未收入的十个原编号题组。历史选择题保留当年课程的答案口径；后来的精测数值和教科书细节不能倒填进旧题。图题改以坐标与运动方向写出完整条件，不复制受限教材的图页。标为“官方”的结论来自公开复习解答，其间的推导与文字为本卷重写。

\src{8.286 Quiz 1 Review Problem 1(a)--(g), 原题 2000 Quiz 1 Q1；公开官方复习解答。}
\begin{exercise}
(a) Doppler 预言声、光的频移是在 xx42 年，求 xx。
(b) 从银河系形状解释晴夜的银河亮带。
(c) Hubble 1929 年最初论文的星系样本，在 $2,18,180,1800$ 四项中选课程所给数。
(d) 均匀 Hubble 流示意图沿直线等距排列 Z、A、B、C、D 五个星系。在 A 看 Z、B、C、D 的速度分别为 $-v_0,+v_0,+2v_0,+3v_0$。分别改以 B 和 C 为参照，给五者速度的方向与倍数，并说明为何不产生宇宙中心。
(e) 宇宙线性尺度增大 $1\%$ 时，今天所定义的粒子视界实际距离增幅是大于、等于还是小于 $1\%$？
(f) 1964 年发现背景辐射的两人及其机构是什么？
(g) 复合后，平均光子间距、典型波长、光子数密度、辐射能量密度和温度各如何随线性尺度 $a$ 变化？
\end{exercise}
\begin{solution}
(a) $1842$ 年，故 xx$=18$。(b) 太阳位于多恒星的扁平盘内，沿盘面看去穿过更多恒星，形成亮带。(c) 官方选择 $18$；官方解答同时指出 Hubble 原图的黑点似有约 $24$ 个，故这里报告的是\emph{课程读本的答案}，不以题选项抹掉史料歧义。
(d) 取向右为正、等距 $L$，$v_0=HL$。B 参照的 Z、A、B、C、D 速度为 $(-2,-1,0,+1,+2)v_0$；C 参照则为 $(-3,-2,-1,0,+1)v_0$。任意换原点都满足 $v_i-v_j=H(x_i-x_j)$，所以观测者均见相同形式的退行规律。(e) $D_p(t)=a(t)\int_0^t c\,\dd t'/a(t')$；后一个共动积分随时间严格增加，故 $D_p$ 增幅大于 $a$ 的 $1\%$。(f) Arno Penzias 与 Robert Wilson，贝尔电话实验室。(g) 依次为 $a,\ a,\ a^{-3},\ a^{-4},\ a^{-1}$。能量密度多出的一个逆幂来自单光子红移 $E=hc/\lambda\propto a^{-1}$。
\end{solution}

\src{8.286 Quiz 1 Review Problem 3(a)--(e), 原题 2007 Quiz 1 Q1；公开官方复习解答。}
\begin{exercise}
(a) 写出 1940 年代稳态理论三位提倡者中至少一人的姓。原卷计分：至少一人正确得 5 分；另两人各加 1 分，每错报一人扣 1 分。
(b) de Sitter 在 1917 年的成就是哪项：测得银河直径十亿光年、用仙女座造父变星证实河外星系、编 103 个星云星团的彗星排除表、提出看似静止而有距离红移的相对论宇宙模型、发现行星周期与轨道半长轴的 $3/2$ 幂关系？
(c) 从贝尔实验室、MIT、普林斯顿、鸽子、旱獭、Hubble 常数、液氦、$7.35\,\mathrm{cm}$ 八项中选\emph{未参与}背景辐射发现故事的两项。原卷评分：只答对一项得 3 分，两项均对得 5 分；每错选一项扣 2 分，最低 0 分。
(d) 恒星光行差显示地球公转速度按预期变化，Foucault 摆显示地球自转。两项发现分别何时完成？从以下四项选择：
(i) 哥白尼在世时；(ii) 他去世后约二十年、三十年；
(iii) 他去世后约一百年；(iv) 他去世后分别约两个世纪、三个世纪。
(e) 在 $1\,\mathrm{AU},100\,\mathrm{kpc},1,10,100,1000\,\mathrm{Mpc}$ 中选课程认为宇宙平均均匀各向同性开始适用的尺度。
\end{exercise}
\begin{solution}
(a) Bondi、Gold、Hoyle。(b) 第四项；de Sitter 当年的坐标描述被看作静止，现代解释承认膨胀。(c) 旱獭和 Hubble 常数。鸽子曾是接收机噪声排查的一部分，普林斯顿团队探寻背景，MIT 也出现在当事人交流背景中；不能仅凭“不是最后发现机构”判错。(d) 选 (iv)：光行差由 Bradley 于 1728 年发现，约在哥白尼去世两世纪后；Foucault 摆实验在 1851 年，约三世纪后。(e) $100\,\mathrm{Mpc}$，是此历史题的粗尺度，并非宣称小于此尺度任何统计量都不均匀。
\end{solution}

\src{8.286 Quiz 1 Review Problem 5(a)--(j), 多届阅读题汇编；公开官方复习解答。}
\begin{exercise}
(a) 在无宇宙常数的课程三分模型中，用何量区分开放、平直、封闭宇宙：微波温度、$\Omega=\rho/\rho_c$、物质/辐射主导还是红移？
(b) 今天 CMB 温度约多少 K？(c) 在附近星系团距离、CMB 观测、大尺度星系聚集、球状星团年龄中，何项支持各向同性？
(d) 仙女座距离在 $10\,\mathrm{kpc}$、$5$ 十亿光年、$2$ 百万光年、$3$ 光年中选何者？
(e) Hubble 定律发表于 1862、1880、1906、1929、1948 中哪年？
(f) Hubble 当年测得 $H^{-1}$ 约 $10^8$、$2\times10^9$、$10^{10}$、$2\times10^{10}$ 年中哪项？
(g) 造父变星怎样用于测距？(h) 旧读本的适用距离在 $10^2,10^4,10^7,10^{10}$ 光年中选哪项？
(i) 写出背景辐射发现者及机构。(j) 当时独立准备探测背景的 Peebles、Dicke、Roll、Wilkinson 在哪一机构？
\end{exercise}
\begin{solution}
(a) 课程的 $\Lambda=0$ 背景模型以 $\Omega\lessgtr1$ 定曲率符号；有暗能量时单靠今天物质密度参数不可照搬。(b) 约 $3\,\mathrm K$，精确口径约 $2.7\,\mathrm K$。(c) CMB 天空各方向近乎一致。(d) 约两百万光年。(e) 1929。(f) 约二十亿年，反映当年的距离标尺误差。(g) 造父变星的周期与平均本征光度相关；测周期求绝对光度，再由观测通量 $F=L/(4\pi D^2)$ 推距离。(h) 旧书题选 $10^7$ 光年；官方后来注明该项随特定读本而来，不宜充当无条件的现代仪器上限。(i) Penzias、Wilson，贝尔电话实验室。(j) 普林斯顿大学。
\end{solution}

\src{8.286 Quiz 1 Review Problem 8(a)--(e), 原题 1996 Quiz 1 Q1；公开官方复习解答。}
\begin{exercise}
(a) Fraunhofer 在 1814--1815 年见到的太阳棱镜谱元素线为亮线还是暗线，为什么？
(b) 仙女座恒星得以分辨并证实为河外系统，发生于 Galileo 1609、Kant 1755、Leavitt 1912、Hubble 1923、Baade/Sandage 1950 年代中的何者？
(c) 判断五句真伪：CN 分子温度曾在 CMB 直接发现二十多年前被测；CN 是将微波反射至地面接收机；CN 给出黑体峰短波侧的等效温度；分析穿过 CN 云的可见星光谱能反推照射它的微波强度；须由云内化学反应才能反推微波温度。
(d) 公元前约 280 年提出日心设想的古希腊思想家是谁？
(e) 静止、透明、无始且恒星均匀分布的宇宙里，Olbers 悖论预言的任一夜空视线表面亮度近太阳表面、月面、无穷大，还是总天空通量约等于太阳？
\end{exercise}
\begin{solution}
(a) 暗线；较冷的太阳外层原子选择吸收连续谱的特定波长。(b) Hubble 于 1923 年观测造父变星并建立距离证据。(c) 依次真、假、真、真、假。CN 激发能级布居改变可见吸收线强度，由此间接推测照射云的背景微波场，而非微波镜面反射；本题与正式 Quiz 1 Q1(c) 的原理相同。(d) Aristarchus；官方注明这是旧读本题，之后学期不宜当公平的必读考点。(e) 每条视线终将止于恒星，\emph{单位立体角表面亮度}近似恒星表面亮度；总亮度并非数学无穷，真实宇宙有限年龄与膨胀使前提失效。
\end{solution}

\src{8.286 Quiz 1 Review Problem 10(a)--(d), 原题 1998 Quiz 1 Q1；公开官方复习解答。}
\begin{exercise}
(a) Einstein 1917 年引进宇宙学常数的初始动机是什么？
(b) de Sitter 1917 年模型当时被认为静止还是膨胀，今天又怎样理解？
(c) 热黑体光子的数密度与能量密度分别正比 $T$ 的几次幂？
(d) 约 $3000\,\mathrm K$ 时物质发生了什么化学/电离态变化，为什么这与后来结构形成有关？
\end{exercise}
\begin{solution}
(a) 使当时相信的静态宇宙在场方程中成为可能。(b) 当年被视作静止，现代由其他坐标辨识其膨胀性质。(c) $n_\gamma\propto T^3,\ \rho_\gamma c^2\propto T^4$，因动量空间热尺度 $\sim\kb T/c$ 提供三个逆长度幂，每光子再贡献一能量幂。(d) 自由电子与核结合成中性原子，即“复合”；光子 Thomson 散射急减，背景光自由传播，普通重子不再被辐射压紧密拖曳，更易向暗物质势阱聚集。小幅暗物质扰动在此以前已可生长，不能说所有结构从此才出现。
\end{solution}

\src{8.286 Quiz 1 Review Problem 14(a)--(e), 原题 2004 Quiz 1 Q1；公开官方复习解答。}
\begin{exercise}
(a) 陈述 Olbers 夜空悖论及主因解法。(b) Planck 单位中 $G$ 为何值；Planck 长在 $10^{-35}\,\mathrm m$、$10^{-15}\,\mathrm m$、$10^{15}\,\mathrm m$、$10^{35}\,\mathrm m$ 中为何项？
(c) 何为宇宙学原理，Hubble 膨胀是否与之相容？
(d) 透明化温度约 $3\times10^a\,\mathrm K$，今日 CMB 约 $3\times10^b\,\mathrm K$，求整数 $a,b$。
(e) 大爆炸后约 $0.01\,\mathrm s$，列出预计占总质量/能量密度超过约 $1\%$ 的粒子种类。
\end{exercise}
\begin{solution}
(a) 若空间透明、恒星恒定且无穷古老，每一方向最终迎上恒星，夜空应像恒星表面一样明亮；宇宙有有限过去与因果视界，宇宙膨胀和恒星演化也使假设失效。(b) 选单位使 $c=\hbar=G=1$，故 $G=1$，长度约 $10^{-35}\,\mathrm m$。(c) 足够大尺度上无特选位置或方向，即统计均匀、各向同性；$v=HD$ 对任意共动观察者都同形，故相容。(d) $a=3,b=0$。(e) 光子、电子、正电子、三味中微子及反中微子；质子、中子虽决定净重子数，未必各占该时刻能量预算的 $1\%$。具体分数依题所采用的热史近似。
\end{solution}

\src{8.286 Quiz 1 Review Problem 15(a)--(b), 原题 2004 Quiz 1 Q2；公开官方复习解答。}
\begin{exercise}
空间站 Alpha-7 静止。一艘射电源飞船以速率 $v_s$ 向站的左方运动，接收者飞船以速率 $v_o$ 向站的右方运动，二者正沿传播方向互相远离，$v_s,v_o<c$。(a) 求接收的频移 $1+z$。(b) 求接收者量得的源退行速度 $v_{\rm tot}$。
\end{exercise}
\begin{solution}
站系中，源向左运动，使沿右方传播的信号在站旁测得 $1+z_s=\sqrt{(1+\beta_s)/(1-\beta_s)}$；右行接收者对该信号再添 $\sqrt{(1+\beta_o)/(1-\beta_o)}$。故
\[
1+z=\sqrt{\frac{(1+\beta_s)(1+\beta_o)}
{(1-\beta_s)(1-\beta_o)}}\,,\qquad
\beta_{s,o}=v_{s,o}/c.
\]
把整体写成 $\sqrt{(1+\beta_{\rm tot})/(1-\beta_{\rm tot})}$ 并平方整理，得
$v_{\rm tot}=(v_s+v_o)/(1+v_sv_o/c^2)$。它小于 $c$，而用普通速度相加会丢掉分母。
\end{solution}

\src{8.286 Quiz 1 Review Problem 16(a)--(f), 原题 2005 Quiz 1 Q1；公开官方复习解答。}
\begin{exercise}
(a) 首个被确立为银河系外的星系是哪一个，怎样测距？
(b) 何为复合；按课程题意星系在其前或后开始形成，为什么？
(c) 高温热气体“配方”所需守恒量在重子数、每粒子能量、质子数、电荷、压强中应选哪些？
(d) 今日 CMB 典型光子的波长在 $2\,\mathrm{fm},2\,\mu\mathrm m,2\,\mathrm{mm},2\,\mathrm m$ 中应选哪项？
(e) 陈述等效原理。(f) 地基仪器为何不易观测 CMB 光谱的短波侧？
\end{exercise}
\begin{solution}
(a) 仙女座 M31；Hubble 辨认造父变星，用周期--光度关系推绝对光度及距离。(b) 核与自由电子结合成中性原子，光子退耦；课程把可见星系的重子坍缩置于复合\emph{后}，因为复合前光子压强妨碍重子落入扰动。暗物质扰动早已能演化，须与“星系形成”口语说法区分。(c) 重子数与电荷；质子数可经弱相互作用变化，压强和单粒子能量由热平衡条件决定。(d) 官方选择 $2\,\mathrm{mm}$，是 CMB 毫米波尺度；按“平均能量 $E$ 转波长 $hc/E$”与按波长谱峰定义会差数倍，故不能把本选择题的粗波长当精确平均值。(e) 小范围内均匀重力与等加速度的实验效应等价；更大范围可由潮汐效应区别。(f) 官方按波长短于约 $0.3\,\mathrm{cm}$ 的地球大气吸收与再散射解释；地面若要进入相关窗口须专门校正，高空气球或卫星更便于观测。
\end{solution}

\src{8.286 Quiz 1 Review Problem 18(a)--(c), 原题 2005 Quiz 1 Q4；公开官方复习解答。}
\begin{exercise}
忽略宇宙膨胀。Xanthu 与 Emmerac 为两艘飞船，信号均为沿 $y$ 方向传播的射电波，取 $0<v<c$。(a) Xanthu 静止于 $(0,a,0)$，Emmerac 沿 $+x$ 轴以 $v$ 运动，穿越原点时收到信号。求 $z$ 的符号和值。(b) Emmerac 静止于原点，Xanthu 沿 $-x$ 方向以 $v$ 运动，在正上方 $(0,a,0)$ 发出垂直信号。再求 $z$。(c) 两图可否视为同一过程的 Lorentz 换系，为什么？发射和接收可取时间间隔的无穷小极限。
\end{exercise}
\begin{solution}
(a) 在 Xanthu 静止系，接收瞬间速度与光方向垂直。相邻波峰在该系的抵达坐标时间隔等于源钟间隔，但移动接收钟少走 $1/\gamma$，因此 $1+z=1/\gamma=\sqrt{1-v^2/c^2}<1$，为蓝移。(b) 在接收者静止系，发射者钟因运动变慢，且发射瞬间速度横向于光路，接收间隔是 $\gamma$ 倍源固有间隔：$1+z=\gamma>1$，为红移。这里“横向”在不同观察系定义，不能交换两式。
(c) 不可。当把(a) Lorentz 变换到 Emmerac 静止系时，光传播方向也因像差取得非零 $x$ 分量；(b)却规定光在 Emmerac 系仍垂直。源、收世界线貌似互换，光的四动量条件不同，故两个真实物理事件组不等价。相同的(a)(b)已分别在本卷 PS 1.2 中出现；此处的(c)是复习卷独立补题。
\end{solution}

\src{8.286 Quiz 1 Review Problem 19(a)--(b), 原题 2007 Quiz 1 Q3；公开官方复习解答。}
\begin{exercise}
双层旋转盘的四辆内车以半径 $R$、切速 $v$ 匀速绕中心，四辆外车固定在同一四条径向臂延长段，以半径 $2R$、切速 $2v<c$ 同步转动。外车编号 5--8，内车编号 1--4，车 1 与车 5 同臂，车 6 位于下一臂。(a) 内车 1 接收四辆外车各自全向辐射的光，求各红移。(b) 一飞船沿经过中心的水平线向右以速率 $u$ 退行，车 6 正位于右方臂端，其瞬时切速度垂直于来自飞船的光路；求车 6 所见红移。
\end{exercise}
\begin{solution}
令 $\gamma_1=(1-v^2/c^2)^{-1/2}$、$\gamma_2=(1-4v^2/c^2)^{-1/2}$。两车相对几何在旋转中固定，连续波峰所需站系传播时间相同；所以站系发射、接收间隔相等。源外车的固有间隔乘 $\gamma_2$ 成站系间隔，内车接收固有间隔再除 $\gamma_1$，四个结果均为
\[
1+z_{5,6,7,8\to1}=\frac{\gamma_2}{\gamma_1}
=\sqrt{\frac{1-v^2/c^2}{1-4v^2/c^2}}.
\]
(b) 可在车 6 旁置一瞬时静止于中心惯性系的虚拟转发器。飞船向右远离该点，入射光向左，给站系纵向因子 $\sqrt{(1+u/c)/(1-u/c)}$；车 6 横向运动，其接收固有时间隔比站系短 $\gamma_2$ 倍。因此
\[
1+z_{S\to6}=\frac1{\gamma_2}
\sqrt{\frac{1+u/c}{1-u/c}}
=\sqrt{\frac{(1-4v^2/c^2)(1+u/c)}{1-u/c}}.
\]
该式可以小于 1：横向接收蓝移可能压过源退行红移。
\end{solution}
